\documentclass[spanish,letterpaper,oneside]{article}
\usepackage[utf8]{inputenc}
\usepackage[T1]{fontenc}
\usepackage[spanish,es-nodecimaldot]{babel}
\usepackage{helvet}
\renewcommand{\familydefault}{\sfdefault}
\usepackage[margin=1.4cm,top=1.1cm,bottom=1.15cm]{geometry}
\usepackage{amsmath,amssymb,array,tabularx,colortbl,graphicx,tikz,fancyhdr}
\usetikzlibrary{babel}
\usepackage{xcolor}

\definecolor{iiiepeNavy}{HTML}{163B5C}
\definecolor{iiiepeBlue}{HTML}{1F6F9F}
\definecolor{iiiepeAqua}{HTML}{20B7A8}
\definecolor{iiiepeOrange}{HTML}{F28C28}
\definecolor{iiiepePale}{HTML}{EEF5F8}
\definecolor{iiiepePaper}{HTML}{F6F8F7}
\definecolor{iiiepeInk}{HTML}{1E2A32}

\setlength{\parindent}{0pt}
\setlength{\parskip}{3pt}
\arrayrulecolor{iiiepeNavy!60}

\pagestyle{fancy}
\fancyhf{}
\renewcommand{\headrulewidth}{0pt}
\fancyfoot[L]{\scriptsize\color{iiiepeNavy} IIIEPE · MAEM · PFM02 Temas Selectos de Matemáticas I}
\fancyfoot[C]{\scriptsize\color{iiiepeNavy} Antología General (SD26)}
\fancyfoot[R]{\scriptsize\color{iiiepeNavy} Pág. \thepage}

\begin{document}
\color{iiiepeInk}

% ==================== PÁGINA 1: PORTADA GENERAL DE LA ANTOLOGÍA ====================
\thispagestyle{empty}
\vspace*{0.35cm}
\begin{center}
  \IfFileExists{../actividad-01-que-sabemos-y-como-lo-ensenamos/img/iiiepe/marca-iiiepe.png}{%
    \includegraphics[height=1.95cm]{../actividad-01-que-sabemos-y-como-lo-ensenamos/img/iiiepe/marca-iiiepe.png}%
  }{%
    \textbf{\Large IIIEPE}%
  }
  \vspace{0.3cm}

  {\large\bfseries Instituto de Investigación, Innovación y Estudios de Posgrado para la Educación del Estado de Nuevo León}\\[4pt]
  {\normalsize Maestría en Aprendizaje y Enseñanza de las Matemáticas (MAEM)}
  \vspace{0.55cm}

  \colorbox{iiiepeNavy}{%
    \parbox{0.94\textwidth}{%
      \centering\vspace{14pt}
      {\LARGE\bfseries\color{white} ANTOLOGÍA GENERAL DE EVIDENCIAS,\\[4pt] RETOS Y EVALUACIONES}\\[8pt]
      {\normalsize\color{iiiepePale} Unidad de Aprendizaje: PFM02 Temas Selectos de Matemáticas I (Bloques I y II · Cohorte SD26)}
      \vspace{14pt}
    }%
  }
\end{center}

\vspace{0.5cm}
\renewcommand{\arraystretch}{1.28}
\noindent\begin{tabularx}{\textwidth}{|>{\bfseries\raggedright\arraybackslash}p{0.28\textwidth}|>{\raggedright\arraybackslash}X|}
\hline
\rowcolor{iiiepePale} Estudiante & Mtro. Martín Jonathan de la Cruz Muñoz \\\hline
Matrícula & 00874 \\\hline
\rowcolor{iiiepePale} Correo institucional & martin.delacruz@campus.iiiepe.edu.mx \\\hline
Docente titular & Mtra. Denis Concepción Escobedo Martínez \\\hline
\rowcolor{iiiepePale} Programa y cohorte & Maestría en Aprendizaje y Enseñanza de las Matemáticas (SD26 · Ago--Nov 2026) \\\hline
Estructura del volumen & Parte I: Portafolio del Bloque I (Sesiones 1--5) · Parte II: Bloque II, Retos presenciales y Parcial 1 (Sesiones 8--13) \\\hline
\rowcolor{iiiepePale} Extensión total & 104 páginas foliadas en formato Carta (vertical uniforme listo para impresión y engargolado) \\\hline
Fecha de integración & 21 de noviembre de 2026 · Monterrey, Nuevo León \\\hline
\end{tabularx}

\vspace{0.55cm}
\noindent\colorbox{iiiepePale}{%
  \parbox{\dimexpr\textwidth-2\fboxsep\relax}{%
    \vspace{4pt}
    {\normalsize\bfseries\color{iiiepeNavy} Propósito y organización de esta Antología}\\[5pt]
    \small Este volumen único reúne de forma secuencial e íntegra la totalidad de las producciones académicas de \textit{Temas Selectos de Matemáticas I}:
    \begin{itemize}\setlength{\itemsep}{3pt}\setlength{\topsep}{3pt}
      \item \textbf{Parte I (págs. 3--41):} Portafolio completo de evidencias del Bloque I (fracciones, decimales, operaciones con signo, proporcionalidad directa/inversa y ecuaciones lineales).
      \item \textbf{Parte II (págs. 42--104):} Reportes completos del Bloque II (Actividades 1 a 6, con rigor formal y bibliografía académica), anteponiendo a cada sesión sus \textbf{hojas de retos presenciales resueltos}, los esquemas geométricos en \textbf{TikZ/GeoGebra}, el \textbf{Examen Parcial 1} resuelto, las \textbf{hojas verdes de la Sesión 12} (Guion de la Habilidad y Solución a los Retos) y el \textbf{Trabajo Final} con referencias digitales a GeoGebra y video en línea.
    \end{itemize}
    \vspace{2pt}
  }%
}

\vfill
\begin{center}
  {\footnotesize\color{iiiepeNavy} Volumen impreso integrado · PFM02 Temas Selectos de Matemáticas I · IIIEPE (SD26)}
\end{center}

% ==================== PÁGINA 2: ÍNDICE MAESTRO PAGINADO (1-104) ====================
\newpage
\begin{center}
  {\large\bfseries\color{iiiepeNavy} Índice Maestro Paginado de la Antología General}\\[2pt]
  {\footnotesize Ubicación exacta por bloque, actividad, hojas de retos presenciales y foliado continuo (Págs. 1--104)}
\end{center}
\vspace{2pt}

\renewcommand{\arraystretch}{1.10}
{\scriptsize
\noindent\begin{tabularx}{\textwidth}{|>{\raggedright\arraybackslash}p{0.23\textwidth}|>{\raggedright\arraybackslash}X|>{\centering\arraybackslash}p{0.08\textwidth}|>{\centering\arraybackslash}p{0.10\textwidth}|}
\hline
\rowcolor{iiiepeNavy}
\color{white}\textbf{Sección / Producto} &
\color{white}\textbf{Contenido y validación (TikZ / GeoGebra / Fuentes)} &
\color{white}\textbf{Ext.} &
\color{white}\textbf{Páginas} \\\hline
Preliminares & Portada General e Índice Maestro Paginado de la Antología & 2 pp. & \textbf{1--2} \\\hline
\rowcolor{iiiepeAqua!22}
\multicolumn{4}{|l|}{\textbf{\color{iiiepeNavy} PARTE I · PORTAFOLIO DE EVIDENCIAS DEL BLOQUE I (SESIONES 1--5)}} \\\hline
Bloque I · Portada & Portada e índice particular del Portafolio Bloque I & 1 p. & 3 \\\hline
\rowcolor{iiiepePale}
Bloque I · Act. 1 & Conversión de fracciones a decimales y viceversa: densidad y secuencias & 9 pp. & 4--12 \\\hline
Bloque I · Act. 2 & Suma y resta con fracciones: equivalencia y recta numérica & 7 pp. & 13--19 \\\hline
\rowcolor{iiiepePale}
Bloque I · Act. 3 & Multiplicación, división y números con signo: modelos geométricos & 10 pp. & 20--29 \\\hline
Bloque I · Act. 4 & Proporcionalidad directa e inversa en un lanzador & 6 pp. & 30--35 \\\hline
\rowcolor{iiiepePale}
Bloque I · Act. 5 & Ecuaciones lineales y propiedades de la igualdad & 6 pp. & 36--41 \\\hline
\rowcolor{iiiepeAqua!22}
\multicolumn{4}{|l|}{\textbf{\color{iiiepeNavy} PARTE II · BLOQUE II: ENSEÑAR DESDE EL DESAFÍO (SESIONES 8--13, RETOS Y PARCIAL 1)}} \\\hline
Separador Parte II & Portadilla divisora del Bloque II y localizador de retos físicos & 1 p. & \textbf{42} \\\hline
\rowcolor{iiiepePale}
Bloque II · Act. 1 & ¿Qué sabemos y cómo lo enseñamos? + \textbf{Retos 1, 2 y 3} (TikZ y parábola $72\text{ m}^2$) + respaldo Canva Equipo 4 & 10 pp. & 43--52 \\\hline
Bloque II · Act. 2 (Retos) & \textbf{Hojas de Retos Resueltos (Sesión 9):} Reto 1 Mediatriz, Reto 2 Bisectriz y Reto 3 Circuncentro (esquemas TikZ) & 3 pp. & \textbf{53--55} \\\hline
\rowcolor{iiiepePale}
Bloque II · Act. 2 (Reporte) & Reporte \textit{El lugar perfecto}: Figura 1 TikZ + 8 construcciones GeoGebra y demostración congruencia LAL/AAL & 8 pp. & 56--63 \\\hline
Bloque II · Act. 3 (Retos) & \textbf{Hojas de Retos Resueltos (Sesión 10):} Reto 1 ($180^\circ$), 2 circunferencias púrpuras ($n=6,8$), ficha 3 polígonos y tabla de 9 polígonos hasta $n=24$ & 4 pp. & \textbf{64--67} \\\hline
\rowcolor{iiiepePale}
Bloque II · Act. 3 (Reporte) & Reporte \textit{Universo 3-60}: triangulación TikZ, demostración por 3 casos de Euclides (ángulo inscrito) y GeoGebra & 9 pp. & 68--76 \\\hline
Examen Parcial 1 & \textbf{Parcial 1. Sospechoso, muy sospechoso:} 5 afirmaciones resueltas (Cierta, Falsa, Cierta, Falsa, Cierta) + 5 esquemas TikZ & 2 pp. & \textbf{77--78} \\\hline
\rowcolor{iiiepePale}
Bloque II · Act. 4 & \textit{Construyendo aprendizaje a través de habilidades}: SEP (1994), Alsina (2023) y esquema TikZ de convergencia & 7 pp. & 79--85 \\\hline
Bloque II · Act. 5 (Hojas) & \textbf{Hojas Verdes (Sesión 12):} \textbf{Pág. 86 = Guion de la Habilidad (\textit{Generalizar})}; \textbf{Pág. 87 = Solución a los 7 retos + \textit{Comunicar}} (TikZ) & 2 pp. & \textbf{86--87} \\\hline
\rowcolor{iiiepePale}
Bloque II · Act. 5 (Reporte) & Reporte \textit{Hacer visible el desarrollo de habilidades}: cuadrícula $2\times 4$ TikZ (pág. 90) y validación GeoGebra ($32\text{ N}$) & 9 pp. & 88--96 \\\hline
Bloque II · Act. 6 (Final) & \textbf{Trabajo Final:} \textit{Una habilidad, un PDA, cinco minutos} (prensa hidráulica $8\text{ N}\to 32\text{ N}$, $45\,271\text{ Pa}$, $W=0.32\text{ J}$; video \texttt{youtu.be/xZOLQRSTyQE}) & 8 pp. & 97--104 \\\hline
\end{tabularx}
}

\vspace{4pt}
\noindent\colorbox{iiiepePale}{%
  \parbox{\dimexpr\textwidth-2\fboxsep\relax}{%
    \scriptsize\textbf{Nota para impresión y encuadernado:} Las 9 páginas diseñadas originalmente en formato apaisado (\textbf{págs. 53--55, 64--67 y 86--87}) fueron normalizadas a Carta vertical respetando su escala y orientación de lectura para encuadernación uniforme sin recortes ni saltos de margen. Los archivos complementarios digitales (\texttt{.ggb} de GeoGebra y video \texttt{.mp4}/YouTube) acompañan este tomo en la memoria USB.
  }%
}

% ==================== PÁGINA 3 (PÁG. 42 GLOBAL): SEPARADOR PARTE II ====================
\newpage
\vspace*{2.2cm}
\begin{center}
  \colorbox{iiiepeNavy}{%
    \parbox{0.88\textwidth}{%
      \centering\vspace{16pt}
      {\Large\bfseries\color{iiiepeAqua} PARTE II · BLOQUE II}\\[6pt]
      {\Huge\bfseries\color{white} Enseñar desde el desafío}\\[6pt]
      {\normalsize\color{iiiepePale} Sesiones 8 a 13 · Reportes Académicos, Hojas de Retos Resueltas,\\[2pt] Guion de la Habilidad, Examen Parcial 1 y Trabajo Final}
      \vspace{16pt}
    }%
  }
\end{center}

\vspace{0.8cm}
\noindent\begin{tabularx}{\textwidth}{|>{\bfseries\raggedright\arraybackslash}p{0.26\textwidth}|>{\raggedright\arraybackslash}X|>{\centering\arraybackslash}p{0.14\textwidth}|}
\hline
\rowcolor{iiiepeNavy}
\color{white}\textbf{Componente en Parte II} &
\color{white}\textbf{Localización directa en este tomo} &
\color{white}\textbf{Páginas} \\\hline
\rowcolor{iiiepePale}
Actividad 1 (Sesión 8) & Reporte completo y Retos 1, 2 y 3 con esquemas TikZ + evidencia Canva Equipo 4 & 43--52 \\\hline
Actividad 2 (Sesión 9) & Hojas de retos (págs. 53--55) + Reporte \textit{El lugar perfecto} y GeoGebra (págs. 56--63) & 53--63 \\\hline
\rowcolor{iiiepePale}
Actividad 3 (Sesión 10) & 4 hojas de retos púrpuras (págs. 64--67) + Reporte \textit{Universo 3-60} (págs. 68--76) & 64--76 \\\hline
Examen Parcial 1 & \textit{Parcial 1. Sospechoso, muy sospechoso} (tabla de 5 afirmaciones + 5 TikZ) & 77--78 \\\hline
\rowcolor{iiiepePale}
Actividad 4 (Sesión 11) & Reporte \textit{Construyendo aprendizaje a través de habilidades} (SEP 1994 y Alsina 2023) & 79--85 \\\hline
Actividad 5 (Sesión 12) & Hojas verdes: Pág. 86 \textit{Guion de Habilidad}; Pág. 87 \textit{Retos} + Reporte (págs. 88--96) & 86--96 \\\hline
\rowcolor{iiiepePale}
Actividad 6 (Sesión 13) & Reporte de Trabajo Final: \textit{Una habilidad, un PDA, cinco minutos} ($32\text{ N}$ y $45\,271\text{ Pa}$) & 97--104 \\\hline
\end{tabularx}

\vspace{0.6cm}
\noindent\colorbox{iiiepePale}{%
  \parbox{\dimexpr\textwidth-2\fboxsep\relax}{%
    \small\textbf{Localizador rápido de retos presenciales y examen:}
    \begin{itemize}\setlength{\itemsep}{2pt}\setlength{\topsep}{2pt}
      \item \textbf{Págs. 44--47:} Retos de Actividad 1 (área/perímetro y optimización $72\text{ m}^2$).
      \item \textbf{Págs. 53--55:} Retos 1, 2 y 3 de Actividad 2 (mediatriz, bisectriz y circuncentro).
      \item \textbf{Págs. 64--67:} Retos de Actividad 3 ($180^\circ$, 2 circunferencias, ficha 3 polígonos y tabla de 9 polígonos).
      \item \textbf{Págs. 77--78:} Examen Parcial 1 completo y justificado.
      \item \textbf{Págs. 86--87:} Hojas verdes de Actividad 5 (Guion de la Habilidad + Solución a los retos de Sesión 12).
    \end{itemize}
  }%
}

\end{document}
